\chapter{Den demokratiske myten}

\begin{motto}
Det er lett å kalle ein smak demokratisk\\ når han alt har vunne. Vanskelegare er det\\ å seie kva han ville gjort om han tapte.
\end{motto}

Det finst ein periode i designhistoria som faget framleis vender seg mot når det vil minne seg sjølv om at det ein gong stod for noko godt. Det er den skandinaviske. Mellom mellomkrigstida og dei tidlege syttiåra voks det i Norden fram ein måte å lage ting på, og ein måte å snakke om det på, som skil seg frå alt anna i denne boka ved éin ting: han hevda ikkje berre å vere vakker eller rasjonell, men \emph{god}. Ikkje god i tydinga velforma, men god i moralsk tyding — sosial, demokratisk, til gagn for dei mange. Lyse tre, mjuke former, reine flater, ærlege materiale; ein stol kvar arbeidar kunne eige, ein kopp som var like rett i ei arbeidarstove som i ein direktørheim. Her, vil sigerssoga seie, fann designfaget endeleg ikkje berre ein stil, men ein \emph{etikk} — eit grunnlag i det demokratiske idealet. Dette kapitlet skal fortelje korleis det vart sagt, og deretter spørje kva som eigenleg vart lova, og kva som vart levert. For påstanden var stor, og som vi skal sjå: han vart aldri sett på prøve av dei som bar han fram.

Før vi går til historia, eit ord om kva slags påstand vi har for oss, for han er av eit anna slag enn dei føregåande kapitla. Då den amerikanske stylisten sa at ein bil var god, meinte han at han selde; det var ein ærleg, om enn kynisk, dom. Då Rams seier at ein radio er godt design, meiner han at han er sakssvarande og varig; det er ein dom med kriterium ein kan diskutere. Men då den nordiske formgjevaren seier at ein stol er god, meiner han noko som glir mellom fingrane i det du grip etter det: at stolen er vakker, ja, men vakker \emph{på ein måte som er moralsk}, fordi han er enkel, ekte, til folket. Det er ein dom som hentar autoriteten sin frå estetikken og bruksretten sin frå etikken, og som aldri seier kvar den eine sluttar og den andre tek til. Heile dette kapitlet er eit forsøk på å halde dei to frå kvarandre lenge nok til å sjå at faget aldri gjorde det sjølv — og at det var nett i den uklårleiken det fann sin mest varige sjølvros.

Lat oss byrje med orda, for orda kom før tinga. Allereie i 1919 skreiv den svenske kunsthistorikaren Gregor Paulsson eit lite skrift med ein tittel som skulle bli eit program for eit halvt hundreår: \textit{Vackrare vardagsvara} — vakrare kvardagsvarer. Tanken var enkel og forførande. Industrien produserer uansett; spørsmålet er berre om han produserer stygt eller vakkert, og kunstnaren sin plikt er å gå inn i fabrikken og sjå til at det vakre vinn. Slöjdföreningen i Sverige gjorde dette til oppdrag. Bak låg ein optimisme av eit slag faget sjeldan har sett sidan: trua på at smak kunne folkeopplysast, at det vakre kunne demokratiserast som lese- og skrivekunsten, at ein heil folkesetnad kunne lærast opp til å ville det rette. Det er ein vakker tanke. Men legg merke til kva han allereie inneheld og ikkje seier: han føreset at nokon veit kva den vakre kvardagsvara er, før folkeopplysinga byrjar. Folket skal lærast opp til å ville henne. Den som lærer opp, må då alt kjenne svaret. Heile prosjektet kviler på at det finst eit fasitsvar om form som læraren har og eleven manglar — og det svaret vart aldri skrive ned, fordi det aldri var ein teori. Det var ein smak. Ein bestemt borgarleg, nordisk, modernistisk smak, som no skulle delast ut til folket som om det var kunnskap.

\section{Stockholm 1930 og den vakre kvardagen}

Tinga kom for alvor i 1930. Stockholmsutstillinga det året, teikna i hovudsak av Gunnar Asplund, var Nordens funksjonalistiske gjennombrot — lette paviljongar i glas og stål, ein arkitektur som proklamerte brot med det tunge og det historiske. Året etter kom manifestet \textit{acceptera}, signert av Asplund, Paulsson og fleire, eit kampskrift for det moderne mot det sentimentale og bakoverskodande. Funksjonalismen kom altså til Norden, slik han kom til Tyskland og elles i Europa, som ein internasjonal modernisme; det var ikkje noko spesifikt nordisk ved han då. Det nordiske kom etterpå, og det kom som ei \emph{mjuking} av funksjonalismen — der den tyske og den franske greina dyrka stål, glas og det maskinelle, dyrka den nordiske etter kvart treet, det handlaga, det organiske, det som ein del seinare kalla ein meir human modernisme.

Her må vi vere historisk presise, for sjølvbiletet er upresist. Funksjonalismen i Norden var aldri eitt prosjekt. Han var fleire, og dei var usamde. I Sverige drog Paulsson og Slöjdföreningen i sosialdemokratisk lei: det handla om bustadpolitikk, om standardisering, om varer til dei mange. I Danmark gjekk utviklinga ein heilt annan veg, gjennom møbelsnikkarlauget og dei årlege snikkarutstillingane, der det avgjerande var handverket, treet og den foredla einskildgjenstanden — Kaare Klint sin tolmodige reform av møbeltypane, ikkje massevara. I Finland var Alvar Aalto noko for seg sjølv, like mykje arkitekt og kunstnar som reformator. Og Noreg var, som så ofte, i utkanten av dei tre, utan ein Aalto og utan eit dansk snikkarlaug. Det fanst, kort sagt, ikkje ein skandinavisk designteori i mellomkrigstida. Det fanst tre–fire nasjonale rørsler med ulike mål, ulike materialpreferansar og ulike syn på tilhøvet mellom handverk og industri. Det dei seinare skulle få felles, var ikkje ein tanke. Det var eit namn, og namnet vart funne opp i utlandet.

\section{Eksportmyten}

For det avgjerande grepet i dette kapitlet — og det vi skuldar Kjetil Fallan og dei forskarane han samla i \textit{Scandinavian Design: Alternative Histories} — er å sjå at sjølve omgrepet «Scandinavian Design» ikkje er ein teori frå Norden, men ei \emph{merkevare} laga for ein marknad\autocite{fallan2012scandinavian}. Han vart til i etterkrigstida, og han vart til utanfor Norden, retta utover. Den store vandreutstillinga som heitte nett \textit{Design in Scandinavia}, og som turnerte USA og Canada frå 1954 til 1957, er det tydelegaste dømet: ei sams nordisk framsyning, sett saman ikkje for nordiske auge, men for det amerikanske, finansiert med eit blikk på den amerikanske kjøpekrafta og den amerikanske lengten etter eit alternativ både til den kalde europeiske modernismen og til den amerikanske stylingen vi såg i førre kapittel. Den britiske marknaden vart opna parallelt, gjennom importørar som Finmar, som selde nordiske møblar inn i eit etterkrigs-Storbritannia som svalt etter lyst tre og reine liner\autocite{davies1998}. «Scandinavian Design» var, med andre ord, frå fyrste stund eit eksportomgrep — ein måte å selje nordiske varer til folk som ikkje budde i Norden, ved å gje dei ei felles, gjenkjenneleg identitet dei ikkje hadde heime.

Og merkevara verka, fordi ho fortalde ei historie kjøparen ville høyre. Ho fortalde at desse tinga ikkje berre var smakfulle, men \emph{dydige}: laga i små, like, sosialdemokratiske velferdssamfunn langt mot nord, av ærlege materiale, av handverkarar som ikkje var korrumperte av kommersialismen, til folk og ikkje til profitt. Dette var ei forteljing like mykje som det var ein stil, og dei to gjorde kvarandre sterke: det vakre vart prov på det gode, og det gode adla det vakre. Halén og Wickman, som femti år seinare gjekk attende til materialet for å sjå kva som var sant, kalla boka si \textit{Scandinavian Design Beyond the Myth} — av di det er nett ein myte dei finn\autocite{halen2003}. Ikkje ei løgn; ein myte, i den presise tydinga: ei forteljing som gjorde reelt arbeid, som batt saman ulike ting til éin meiningsfull figur, og som vart trudd fordi ho var nyttig og fager, ikkje fordi ho var prøvd. Bak figuren låg det reelle og dugande objekt, store formgjevarar, ekte handverk. Men figuren sjølv — eininga, teorien, den sams etikken — var bygd i ettertid og utanfrå, for å selje.

\section{Det demokratiske som vart aldri prøvt}

Lat oss då ta det demokratiske idealet på alvor, alvorlegare enn faget sjølv tok det, og spørje korleis ein ville visst om det heldt. Påstanden var at god form skulle vere til gagn for dei mange — tilgjengeleg, rimeleg, frigjerande. Det er ein påstand med eit innhald, og ein påstand med eit innhald kan i prinsippet prøvast. Vart dei vakre kvardagsvarene faktisk kvardagsvarer? Eigde arbeidaren stolen? Var den nordiske designen rimeleg for dei mange, eller var han, då som no, ein smak for den utdanna middelklassen som hadde råd til å betale for det enkle? Vi veit svaret, og svaret er ubehageleg: mykje av det som vert feira som demokratisk design var, og er, dyrt. Ein Wegner-stol, ein Jacobsen-stol, eit Aalto-glas — dette er statusvarer for ein velståande smak, ikkje massevarer for arbeidaren. Det demokratiske låg i \emph{retorikken} om objektet, ikkje i objektet sin faktiske plass i samfunnet.

Men sjå no kva som er det eigentlege poenget for denne boka, og det er ikkje at idealet svikta i praksis — alle ideal sviktar noko i praksis. Poenget er at faget aldri \emph{sjekka}. Det fanst inga undersøking, ingen vilje til å falsifisere, ingen som spurde: held påstanden vår om at vi lagar demokratisk design, mot tala for kven som faktisk eig tinga? Eit fag med eit fundament ville teke sin eigen sentralpåstand og prøvd å skyte han ned, fordi det er slik kunnskap veks. Designfaget tok i staden den demokratiske påstanden og gjorde han til ein \emph{sjølvskildring} som ikkje trong prøvast, fordi han kjentest sann og smigra utøvaren. Det demokratiske idealet vart aldri operasjonalisert — aldri omsett til noko ein kunne måle — og difor heller aldri falsifisert. Det levde, og lever, som ein dyd faget tileignar seg gratis, utan å betale prisen ein dyd kostar: risikoen for å bli motbevist.

Det er verdt å dvele ved kor lett det hadde vore å gjere noko anna. Å operasjonalisere idealet krev ikkje meir enn å spørje kva «demokratisk» skulle bety, og så sjå etter. Tyder det at tinget er rimeleg? Då kunne ein samanlikne pris mot inntekt, og sjå at den feira stolen kosta ei månadsløn. Tyder det at det er produsert for dei mange? Då kunne ein telje opplag, og sjå at mykje av kanonen vart laga i små seriar for ein smal marknad. Tyder det at det frigjer brukaren, gjer livet betre på ein måte ein kan peike på? Då måtte ein definere kva betre liv er, og det er nett den teorien om form og menneske faget aldri eigde. Kvar av desse tolkingane er prøvbar; ingen av dei vart prøvd. I staden vart «demokratisk» verande eit ord ein heng på tinget, ikkje ein eigenskap ved tinget — meir lik ein heider enn ein hypotese. Og ein heider treng ein ikkje sjekke. Han treng berre delast ut, helst til seg sjølv. Den som vil sjå skilnaden mellom eit fag og ei rørsle, kan sjå han her: rørsla feirar sine verdiar, faget prøver å motbevise sine, og den nordiske tradisjonen valde å feire.

Dette har ein klangbotn også her heime, og han er verd å nemne, for boka er skriven inn i nett den arven. Den norske designtradisjonen, og det miljøet ho høyrer til, har teke imot den demokratiske sjølvforståinga utan å arve verken Aalto eller det danske snikkarlauget — vi fekk forteljinga om den gode, sosiale, nøkterne nordiske formgjevinga, og gjorde henne til vår eigen utan å ha bygd henne. På ein skule som denne lever framleis overtydinga om at det reine, det nedtona og det sosialt ansvarlege er same sak, og at den som meistrar formspråket difor òg eig etikken. Det er ein arv som kjennest god å bere, og det er det som gjer han farleg. For den som ber han, har fått ein moralsk autoritet utan å ha bygd grunnlaget under han, og merkar det ikkje før nokon ber honom grunngje kvifor det enkle skulle vere det rette — og då har han berre forteljinga att.

\vignett

Og dette er det djupaste i kapitlet. Då den nordiske formgjevaren sa at ein stol var god, kva sa han eigentleg? Han sa to ting på éin gong, og blanda dei med vilje. Han sa at stolen var velforma — ein estetisk dom — og han sa at stolen var \emph{moralsk betre} fordi han var enkel, ærleg, ikkje-prangande. Men det andre følgjer ikkje av det fyrste. At ei form er rein, gjer henne ikkje god mot menneske; at ho er pyntelaus, gjer henne ikkje demokratisk; reinleik kan like gjerne vere asketisk overklassesmak som folkeleg dyd. Den nordiske tradisjonen sin store bragd var retorisk: han fekk ein bestemt borgarleg smak for det enkle til å framstå som ein etikk, og han gjorde det så fullkome at vi enno trur på det. Men ein smak heva til moral utan grunnlag er ikkje ein etikk. Det er ein smak med god samvitsattest. Og ein attest faget skreiv ut til seg sjølv.

For å vere rettferdig: dette betyr ikkje at den nordiske designen var dårleg, eller at dei sosiale ambisjonane var uærlege. Mange av dei som bar dei, meinte dei djupt, og noko vart faktisk betre — bustadstandard, kvardagsting, eit visst nivå under det stygge. Innvendinga er ikkje moralsk. Ho er fagleg: at her, der faget kom nærast å ha ein etikk, hadde det ikkje ein etikk, men ei stemning av etikk; ikkje ein teori om kva godt design er for menneske, men ein vakker forteljing om at den nordiske smaken alt var det. Og ein forteljing kan ikkje skilje seg sjølv frå sjølvros, fordi ho manglar reiskapen til å bli motsagd.

\section{Lunning-prisen og den amerikanske konstruksjonen}

For å sjå kor medvite merkevara vart bygd, og kor lite ho hadde med ein nordisk teori å gjere, lønner det seg å følgje pengane og prisane. Det tydelegaste enkeltgrepet er Lunning-prisen, oppretta i 1951 av Frederik Lunning — ikkje ein nordisk kulturinstitusjon, men forretningsmannen som dreiv Georg Jensen sin butikk på Femte Avenue i New York. Prisen gjekk kvart år til to unge nordiske formgjevarar, og over dei tjue åra han fanst, fram til 1970, fungerte han som ein systematisk produsent av nordiske designstjerner for ein amerikansk marknad. Han skapte ein kanon: Tapio Wirkkala, Hans Wegner, Henning Koppel, Grete Prytz Kittelsen frå Noreg. Sjå kva som hender her. Ein amerikansk kjøpmann med ein nordisk butikk i New York oppfann ein pris som gjorde nordiske formgjevarar til namn, slik at varene deira kunne seljast til amerikanarar som dyre kunstgjenstandar. Det er ikkje ein etikk som spreier seg utover frå Norden. Det er ein marknad som byggjer seg ein leverandørmytologi.

Mønsteret går att overalt der ein ser etter det. Vandreutstillinga \textit{Design in Scandinavia} (1954–57) var finansiert med blikk på amerikansk kjøpekraft, slik vi alt har sett. Den store nordiske heimemarknadsmønstringa, H55 i Helsingborg sommaren 1955, var rett nok retta innover, men ho var like fullt ei \emph{utstilling} — ei iscenesetjing av det gode livet i lyse, opne rom, ikkje ei dokumentasjon av korleis folk faktisk budde. Kjetil Fallan og medforfattarane hans i \textit{Scandinavian Design: Alternative Histories} arbeider nett med å skilje denne iscenesetjinga frå den faktiske produksjons- og forbrukshistoria, og det dei finn, er gong på gong at det feira utvalet er eit smalt toppsjikt løfta fram som om det representerte heile\autocite{fallan2012scandinavian}. Den breie nordiske kvardagsproduksjonen — det meste folk faktisk kjøpte og brukte — var verken særleg nordisk i uttrykk eller særleg vakkert; det var billege, importerte, etterlikna varer som ikkje passar inn i forteljinga, og som forteljinga difor stilt let liggje. Lees-Maffei og Fallan har seinare synt korleis nett denne mekanismen — at ein nasjonal designidentitet vert komponert i ettertid for ein global marknad — er regelen og ikkje unntaket i moderne designhistorie\autocite{fallan2016designing}.

\section{Prøven faget aldri tok: prisen mot løna}

Lat oss så gjere det faget aldri gjorde, og prøve den demokratiske påstanden mot tal. Påstanden var at dette var form for dei mange — rimeleg, tilgjengeleg, folkeleg. Det er ein prøvbar påstand, og prøven er enkel: kva kosta tinga, målt mot kva folk tente? Ta nokre av kanonobjekta. Hans Wegner sin runde stol frå 1949 — den som amerikanarane rett og slett kalla «The Chair» etter at Kennedy og Nixon sat i han under TV-debatten i 1960 — var handbygd i edeltre hjå møbelsnikkaren Johannes Hansen og kosta deretter. Arne Jacobsen sin Egg-stol og Svane-stol frå 1958, utvikla for SAS Royal Hotel i København, var og er luksusvarer, produserte av Fritz Hansen til prisar som la dei langt utanfor rekkjevidde for arbeidaren dei retorisk vart laga på vegner av. Eit Aalto-glas, ein Jacobsen-stol, eit Georg Jensen-fat: dette er, då som no, statussymbol for ein velståande, utdanna smak.

Set dette opp mot retorikken, og spriket er pinleg. «Vakrare kvardagsvarer» skulle løfte smaken til heile folket; mykje av det som faktisk vart kanonisert under det namnet, var einskildgjenstandar i edle materiale, produserte i små seriar, prisa for eit toppsjikt. Det fanst rett nok ei genuint demokratisk line — den svenske bustadpolitikken, kooperasjonen, standardiseringa, IKEA frå 1943 og 1958 som faktisk sette rimelege møblar i mange heimar — men nett denne lina vert som regel \emph{ikkje} rekna med i den feira «skandinaviske designen»; ho er for billeg, for masseprodusert, for lite kunstnarleg til å pryde forteljinga. Kanonen valde Wegner framfor IKEA, det dyre handverket framfor den rimelege flatpakka, og kalla deretter heile greia demokratisk. Den verkeleg folkelege produksjonen vart skriven ut av soga om den folkelege designen.

Men det er ikkje sjølve spriket som er bokas poeng — alle ideal sviktar noko i møte med marknaden. Poenget er, som før, at faget aldri sjekka. Det gjorde inga undersøking av kven som faktisk eigde dei feira tinga; det samanlikna ingen pris med noka løn; det spurde aldri om sjølvbiletet heldt. Det hadde vore lett. Ein kunne ha teke ti kanonobjekt, funne utsalsprisen, delt på ei industriarbeidarløn og lese av talet. Resultatet ville ha vore øydeleggjande for forteljinga, og det er truleg difor rekninga aldri vart gjord. Eit fag med eit fundament prøver å skyte ned sin eigen sentralpåstand, fordi det er slik kunnskap veks; den nordiske tradisjonen gjorde påstanden om demokrati til ein heider han delte ut til seg sjølv, og ein heider treng ein ikkje sjekke. Halén og Wickman gjekk femti år seinare attende til materialet og fann nett dette — at det demokratiske låg i talemåten om objektet, ikkje i objektet sin plass i samfunnet — og difor kalla dei boka si \textit{Beyond the Myth}\autocite{halen2003}.

\begin{kasus}{«The Chair»: ein folkeleg stol til ei månadsløn}
Hans Wegner sin Rundstol, modell JH501 frå 1949, er det perfekte prøvestykket, fordi han ber heile motsetnaden i seg. Retorisk er han alt den nordiske forteljinga lovar: enkel, ærleg, organisk, reinsa for prål, eit demokratisk svar på den overlessa borgarstolen. Materielt er han ein handbygd luksusgjenstand. Han vart laga i små seriar hjå møbelsnikkaren Johannes Hansen, kvar stol forma og samanføydd for hand i teak eller eik med flettesete, og prisa deretter. Då han i 1950 vart presentert for det amerikanske publikummet og snart kjend rett og slett som «The Chair», var det som eit objekt for kjennarar, ikkje for arbeidarar. Spør den enkle prøven: kor mange dagsverk kosta han for den arbeidaren han skulle frigjere? Svaret — veker, ikkje timar — avgjer saka. Stolen er framifrå. Han er berre ikkje demokratisk i nokon annan forstand enn at han ser enkel ut, og «ser enkel ut» er ein estetisk eigenskap, ikkje ein sosial. Det var nett denne forvekslinga, av eit reint uttrykk med ein folkeleg realitet, heile myten levde av.
\end{kasus}

\laeringsmal{\begin{itemize}
\item Skilje dei tre–fire nasjonale nordiske rørslene (svensk bustadpolitikk, dansk snikkarlaug, finsk Aalto, norsk utkant) og vise at det ikkje fanst éin felles skandinavisk designteori i mellomkrigstida.
\item Gjere greie for korleis omgrepet «Scandinavian Design» vart konstruert som ei eksportmerkevare, gjennom Design in Scandinavia, Finmar og Lunning-prisen.
\item Operasjonalisere den demokratiske påstanden og prøve han empirisk ved å setje pris opp mot løn for kanonobjekt.
\item Drøfte skiljet mellom ein estetisk og ein moralsk dom, og forklare korleis tradisjonen heva ein smak for det enkle til ein etikk utan å byggje grunnlaget under han.
\end{itemize}}

\ovingar{\begin{enumerate}
\item Vel fem feira «skandinaviske» objekt, finn utsalsprisen (historisk eller notid), og rekn han om til talet på dagsverk for ei industriarbeidarløn. Kva seier resultatet om den demokratiske påstanden?
\item Samanlikn eit kanonobjekt (Wegner, Jacobsen, Aalto) med eit samtidig IKEA-produkt. Kvifor reknar forteljinga det eine som «design» og det andre ikkje?
\item Skriv om eit produkt du sjølv kallar «reint» eller «nordisk», og prøv å skilje den estetiske dommen frå den moralske. Kva blir att av «god» når dei to er skilde?
\item Diskuter i seminar: Kvifor sjekka aldri faget om den demokratiske sjølvskildringa heldt? Kva ville det ha kosta tradisjonen å gjere rekninga?
\end{enumerate}}

\vignett

Den skandinaviske myten lova faget det det mest av alt mangla: eit grunnlag som ikkje berre var estetisk, men moralsk — ein grunn til at god form var god \emph{for nokon}, og ikkje berre fin å sjå på. Det er den mest tiltrekkjande av alle dei lovnadene denne boka går gjennom, fordi ho rører ved det rette: at form har følgjer for menneske, og at eit fag om form difor må ha noko å seie om dei følgjene. Men lovnaden vart innfridd med ei merkevare. Faget fekk eit rykte for å vere demokratisk i staden for ein teori om kva demokratisk form er, og det tok ryktet som om det var teorien. Lengre sør hadde nokon mindre tru på stemningar og meir på reglar. Dei skreiv ned ti av dei, og kalla det godt design.

\vidarelesing{\fullcite{fallan2012scandinavian}; \fullcite{halen2003}; \fullcite{davies1998}; \fullcite{fallan2010}; \fullcite{woodham1997}; \fullcite{forty1986}}
